\label{ch:gravity}

\section{Cadena constitutiva del sector gravitatorio: escala cronológica \(108\), reducción tensorial, holonomía y evaluación metrológica de \(G\)}

La gravedad se organiza aqu\'i en cuatro estratos que no deben fundirse:

\begin{equation}
\label{eq:gravity-causal-chain}
\begin{aligned}
\text{estructura cronol\'ogica CT108 y rectas de acci\'on}
&\longrightarrow \text{selector de periodicidad }\theta_{108}^{\rm circ}
\longrightarrow G_{108},\\
\text{incidencia icosa\'edrica}
&\longrightarrow V_{12}\xrightarrow{P_5}V_5
\xrightarrow{\Gamma_5}\operatorname{STF}_3
\xrightarrow{\Lambda(n)}\mathcal H_{\rm TT}(n),\\
\text{ruta holon\'omica enriquecida}
&\dashrightarrow G_{\rm hol},\\
\text{carta metrol\'ogica}
&\longrightarrow 6.67430\times10^{-11}\,
\mathrm{m^3\,kg^{-1}\,s^{-2}}.
\end{aligned}
\end{equation}

La primera l\'inea expresa una condici\'on intr\'inseca y exacta de
periodicidad. La segunda construye
el portador y su reducci\'on transversal sin traza. La tercera conserva una
frontera de construcci\'on expl\'icita. La cuarta es una evaluaci\'on en una
carta decimal declarada: identifica la magnitud, pero no selecciona el car\'acter
gravitatorio.

\section{Recta de magnitud gravitatoria y escala cronol\'ogica \(108\)}

Sean \(c_{\rm int}>0\), \(t_0>0\),
\(\ell_0=c_{\rm int}t_0\) y una secci\'on positiva
\(\hbar_a\) de la recta de acci\'on, con
\(a\in\{\mathrm{pre},\mathrm{ret}\}\). El transductor dimensional es

\begin{equation}
\label{eq:gravity-transducer}
G_a(\theta)
=\theta\frac{c_{\rm int}^{5}t_0^2}{\hbar_a}
=\theta\frac{c_{\rm int}^{3}\ell_0^2}{\hbar_a},
\qquad \theta>0.
\end{equation}

La igualdad de las dos expresiones usa s\'olo
\(\ell_0=c_{\rm int}t_0\), y
\([G]=L^3M^{-1}T^{-2}\). La dimensi\'on queda fijada por la recta; el
problema matem\'atico es seleccionar el car\'acter adimensional \(\theta\).

El calendario CT108 define

\begin{equation}
\label{eq:gravity-clock108}
T_{108}=108t_0,
\qquad
\omega_{108}=\frac{2\pi}{108t_0},
\qquad
R_{108}=\frac{c_{\rm int}}{\omega_{108}}
=\frac{54}{\pi}\ell_0.
\end{equation}

Para cada hoja de acci\'on, sean

\begin{equation}
\label{eq:gravity-clock-energy-mass}
E_{108,a}=\hbar_a\omega_{108},
\qquad
m_{108,a}=\frac{E_{108,a}}{c_{\rm int}^2}.
\end{equation}

Entonces, la longitud de Compton reducida coincide con el radio del ciclo:

\begin{equation}
\label{eq:gravity-compton-clock}
\bar\lambda_{C,a}
=\frac{\hbar_a}{m_{108,a}c_{\rm int}}
=R_{108}.
\end{equation}

\begin{theorem}[Unicidad del selector bajo la condici\'on de periodicidad HMT]
\label{thm:gravity-circular-selector}
Con la convenci\'on reducida
\(r_g=Gm/c_{\rm int}^2\) y la premisa de periodicidad HMT que identifica el radio
gravitatorio del cuanto CT108 con su radio de fase, la condici\'on de cierre
\begin{equation}
\label{eq:gravity-closing-condition}
r_{g,a}(\theta)=R_{108}
\end{equation}
selecciona una \`unica soluci\'on positiva,
\begin{equation}
\label{eq:gravity-theta108}
\boxed{
\theta_{108}^{\rm circ}=\left(\frac{54}{\pi}\right)^2
=295.4525715010569415303525\ldots .}
\end{equation}
En esa realizaci\'on,
\begin{equation}
\label{eq:gravity-four-lengths}
R_{108}=\bar\lambda_{C,a}=r_{g,a}=\ell_P^{\rm circ}.
\end{equation}
\end{theorem}

\begin{proof}
Por \cref{eq:gravity-transducer,eq:gravity-clock-energy-mass},
\[
r_{g,a}(\theta)
=\frac{G_a(\theta)m_{108,a}}{c_{\rm int}^2}
=\theta c_{\rm int}t_0^2\omega_{108}
=\theta\frac{\pi}{54}\ell_0.
\]
Al compararlo con
\(R_{108}=(54/\pi)\ell_0\) se obtiene necesariamente
\(\theta=(54/\pi)^2\). La ecuaci\'on es lineal en \(\theta\) y todos los
factores son positivos, luego la soluci\'on es \`unica.
\end{proof}

La igualdad \(r_g=R_{108}\) no se deduce de la dimensionalidad de
\cref{eq:gravity-transducer}. Es la condición geométrica explícita de esta
realizaci\'on HMT: acci\'on, Compton y retorno peri\'odico deben determinar una misma
longitud. El teorema demuestra que, aceptada esa condición estructural, el
carácter \(\theta\) queda fijado sin usar el valor SI de \(G\).

\begin{remark}[Convenci\'on de radio]
El teorema usa el radio gravitatorio reducido \(Gm/c^2\). Si se sustituye
por el radio de Schwarzschild \(2Gm/c^2\), el selector cambia por un factor
dos. Las dos expresiones corresponden a condiciones geom\'etricas distintas,
por lo que la convenci\'on debe declararse antes del c\'alculo.
\end{remark}

\section{Reciprocidad \(G\leftrightarrow\hbar\) y conservación del área de Planck}

\begin{corollary}[Familia de Planck asociada a la periodicidad temporal de \(108\) fases]
\label{cor:gravity-planck-family}
Para \(G_a=G_a(\theta_{108}^{\rm circ})\),
\begin{align}
\ell_P&=\frac{54}{\pi}\ell_0,
&t_P&=\frac{54}{\pi}t_0,
\label{eq:gravity-planck-length-time}\\
E_{P,a}&=\frac{\hbar_a}{t_P}
=\frac{\pi\hbar_a}{54t_0}
=\frac{h_a}{108t_0},
&\ell_P^2&=\frac{G_a\hbar_a}{c_{\rm int}^3}
=\theta_{108}^{\rm circ}\ell_0^2.
\label{eq:gravity-planck-energy-area}
\end{align}
Adem\'as,
\begin{equation}
\label{eq:gravity-planck-force-power-density}
F_{P,a}=\frac{c_{\rm int}^4}{G_a},
\qquad
P_{P,a}=\frac{c_{\rm int}^5}{G_a},
\qquad
\rho_{P,a}=\frac{\hbar_a}
{(\theta_{108}^{\rm circ})^2c_{\rm int}^5t_0^4}.
\end{equation}
\end{corollary}

\begin{proof}
Las tres primeras identidades proceden de
\(\ell_P=R_{108}\), \(t_P=\ell_P/c_{\rm int}\) y
\(h_a=2\pi\hbar_a\). Para el \`area,
\[
\frac{G_a\hbar_a}{c_{\rm int}^3}
=\theta_{108}^{\rm circ}c_{\rm int}^2t_0^2
=\theta_{108}^{\rm circ}\ell_0^2.
\]
Las expresiones restantes son las composiciones dimensionales de fuerza,
potencia y masa por volumen con la misma secci\'on.
\end{proof}

Sea
\begin{equation}
\label{eq:gravity-Ract}
R_{\rm act}:=\frac{\hbar_{\rm pre}}{\hbar_{\rm ret}}.
\end{equation}
Como la acci\'on aparece en el numerador de la masa cronol\'ogica y en el
denominador de \(G\),
\begin{equation}
\label{eq:gravity-twoway-covariance}
\frac{m_{108,\rm pre}}{m_{108,\rm ret}}=R_{\rm act},
\qquad
\frac{G_{\rm pre}}{G_{\rm ret}}=R_{\rm act}^{-1},
\qquad
G_a\hbar_a
=\theta_{108}^{\rm circ}c_{\rm int}^5t_0^2.
\end{equation}
Por tanto, \(\ell_P^2=G_a\hbar_a/c^3\) no cambia de hoja. La gravedad y
la acci\'on se orientan rec\'iprocamente, mientras el \`area conserva la
geometr\'ia com\'un. Esta cancelaci\'on es la interfaz de confluencia hacia el operador de
\`area del \cref{ch:modular}.
